\documentclass[11pt]{article}
\usepackage{amsmath,amssymb,amsthm,mathtools,hyperref}

\setlength{\parindent}{0pt}
\setlength{\parskip}{0.55em}
\setlength{\textwidth}{6.35in}
\setlength{\oddsidemargin}{0.05in}
\setlength{\evensidemargin}{0.05in}
\setlength{\textheight}{8.75in}
\setlength{\topmargin}{-0.25in}

\newtheorem{theorem}{Theorem}[section]
\newtheorem{lemma}[theorem]{Lemma}
\newtheorem{proposition}[theorem]{Proposition}
\newtheorem{corollary}[theorem]{Corollary}
\theoremstyle{remark}
\newtheorem{remark}[theorem]{Remark}

\newcommand{\Z}{\mathbb{Z}}
\newcommand{\E}{\{-1,0,1\}}
\newcommand{\Rset}{\mathcal{R}}
\newcommand{\Qset}{\mathcal{Q}}
\newcommand{\Dset}{\mathcal{D}}
\newcommand{\res}[2]{\left\langle #1\right\rangle_{#2}}

\title{Round 3 on the proof front: exact critical carry renormalization}
\author{Principia Math Collatz campaign}
\date{August 2026}

\begin{document}
\maketitle

\begin{abstract}
We continue the pinned-carry branch of the zero-trace problem.  The main result is an exact renormalization for the critical change
\[
(k,L)\longmapsto(k+3,L+1).
\]
A pair of realization words has two integer coordinates: its carry and a complementary quotient.  Appending three forward symbols and prepending one backward symbol transforms these coordinates by two coupled affine carry equations.  When the new carry is zero and $L\ge3$, no modular wrap is possible in the first equation.  The new pinned fibre is therefore in bijection with a precisely described subset of the first $64^3$ old carry fibres, selected by one three-step forward test and one one-step backward test.  Each old pair has at most one pinned extension.

We also prove a multi-block no-wrap theorem.  If $J$ critical blocks are stripped from a pinned pair and
\[
(64^3)^J<81^L,
\]
then all $J$ local injections synchronize: the old carry is the base-$64^3$ number formed by the $J$ local forward lift digits, and the complementary quotient has the corresponding base-$81$ expansion.  Applied to a pair of backward length $L_{\rm top}$, this synchronizes every fixed proportion below
\[
\frac{\log 81}{\log(81\cdot64^3)}=0.260472\ldots
\]
of the blocks.  Thus the local injectivity from the earlier framework does multiply for an initial positive fraction of the depth; the remaining obstruction is the first modular wrap.

An explicit example shows why the scalar pinned count is not inductively closed.  A nonzero pair in carry $257792$, rather than carry zero, passes both local tests and produces a pinned pair at $(15,5)$.  Consequently a successful induction must control a decorated low-carry profile, not only $C_{k,L}(0)$.  No proof of the zero-trace conjecture, and no proof of the Collatz conjecture, is claimed.
\end{abstract}

\section{Chosen front and standing input}

\textbf{Chosen front: P3-ZT-PINNED-CARRY, because the smallest plausible proof unit is now the exact synchronization of one critical $3$-forward/$1$-backward block, followed by a determination of where modular wrap first destroys that synchronization.}

Let
\[
F(x)=
\begin{cases}
81x/64,&x\equiv0\pmod{64},\\
(81x+1)/64,&x\equiv15\pmod{64},\\
(81x-1)/64,&x\equiv49\pmod{64},
\end{cases}
\]
with $F$ undefined in the other residue classes.  Let $G$ be its partial inverse,
\[
G(y)=
\begin{cases}
64y/81,&y\equiv0\pmod{81},\\
(64y-1)/81,&y\equiv19\pmod{81},\\
(64y+1)/81,&y\equiv62\pmod{81}.
\end{cases}
\]

For a word
\[
w=(e_0,\ldots,e_{n-1})\in\E^n,
\]
put
\[
N_n(w)=\sum_{j=0}^{n-1}e_j81^{n-1-j}64^j.
\]
Let $R_n(w)\in[0,64^n)$ be its canonical starting value and let $Q_n(w)$ be the endpoint after $n$ steps.  Thus
\[
64^nQ_n(w)=81^nR_n(w)+N_n(w).
\]
Write
\[
\Rset_n=\{R_n(w):w\in\E^n\},
\qquad
\Qset_n=\{Q_n(w):w\in\E^n\}.
\]
For an integer $x$, let $\res{x}{m}$ denote its least residue in $[0,m)$.

\begin{remark}[Results from Rounds 1 and 2 assumed here]
Rounds 1 and 2 proved the existence and uniqueness of $R_n(w)$ and $Q_n(w)$, the injectivity of the two word maps, and the exact identity above.  They also proved the following carry interpretation.  For
\[
A=64^k,
\qquad
M=81^L,
\]
and a pair $(u,v)\in\E^k\times\E^L$, there is a unique carry
\[
c_{k,L}(u,v)\in[0,M)
\]
satisfying
\[
Q_L(v)-R_k(u)\equiv A c_{k,L}(u,v)\pmod M.
\]
If $C_{k,L}(c)$ counts pairs with carry $c$, then
\[
C_{k,L}(0)=N(k,L)+1,
\]
where $N(k,L)$ is the number of nonzero depth-$k$ zero-trace vertices starting a path of length at least $L$.  These statements are assumed in this round.
\end{remark}

\section{The local extension data}

Put
\[
B=64^3.
\]

\begin{lemma}[Perturbation along a fixed word]
\label{lem:perturb}
Let $w\in\E^n$ and let $t\ge0$.  Then the integer
\[
R_n(w)+64^nt
\]
follows the same first $n$ branches as $R_n(w)$, and
\[
F^n\bigl(R_n(w)+64^nt\bigr)=Q_n(w)+81^nt.
\]
\end{lemma}

\begin{proof}
Let $x_j=F^j(R_n(w))$.  We prove by induction that the perturbed orbit has value
\[
x_j+81^j64^{n-j}t
\]
after $j$ steps.  This is true at $j=0$.  For $j<n$, the perturbation term is divisible by $64$, so the two current values have the same residue modulo $64$ and use the same branch error $e_j$.  Therefore
\begin{align*}
F\bigl(x_j+81^j64^{n-j}t\bigr)
&=\frac{81x_j+e_j}{64}+81^{j+1}64^{n-j-1}t\\
&=x_{j+1}+81^{j+1}64^{n-j-1}t.
\end{align*}
At $j=n$ this gives the asserted endpoint.
\end{proof}

Let $u\in\E^k$, let $a\in\E^3$, let $v\in\E^L$, and let $d\in\E$.  Concatenation is written $ua$, while $dv$ means that $d$ is prepended to $v$.

\begin{lemma}[The forward lift digit and the backward lift digit]
\label{lem:lifts}
Define
\[
t_k(u,a)=
\res{\bigl(R_3(a)-Q_k(u)\bigr)81^{-k}}{B}
\]
and
\[
s_L(v,d)=
\res{\bigl(d-64R_L(v)\bigr)64^{-(L+1)}}{81}.
\]
Then
\[
R_{k+3}(ua)=R_k(u)+64^k t_k(u,a)
\]
and
\[
Q_{L+1}(dv)=Q_L(v)+81^L s_L(v,d).
\]
For fixed $u$, the $27$ values $t_k(u,a)$ are distinct.  For fixed $v$, the three values $s_L(v,d)$ are distinct.
\end{lemma}

\begin{proof}
The first $k$ symbols of $ua$ are $u$, so
\[
R_{k+3}(ua)=R_k(u)+64^kt
\]
for a unique $t\in[0,B)$.  Lemma \ref{lem:perturb} shows that the state after those $k$ steps is
\[
Q_k(u)+81^kt.
\]
It follows the final word $a$ precisely when it is congruent to $R_3(a)$ modulo $B$.  Since $81^k$ is invertible modulo $B$, this determines
\[
t=t_k(u,a).
\]
Injectivity of $a\mapsto R_3(a)$ gives distinctness.

Put
\[
p=R_L(v),
\qquad
q=Q_L(v),
\qquad
N=N_L(v),
\qquad
M=81^L.
\]
The affine word identity gives
\[
64^Lq=M p+N.
\]
The numerator of the prepended word is
\[
N_{L+1}(dv)=dM+64N.
\]
Reducing the endpoint congruence for $dv$ modulo $M$ gives
\[
Q_{L+1}(dv)\equiv q\pmod M.
\]
Thus
\[
Q_{L+1}(dv)=q+Ms
\]
for a unique $s\in[0,81)$.  Substituting
\[
N=64^Lq-Mp
\]
into the endpoint congruence modulo $81M$ gives
\begin{align*}
64^{L+1}(q+Ms)
&\equiv dM+64(64^Lq-Mp)\pmod{81M}.
\end{align*}
After cancellation of $64^{L+1}q$ and division by $M$,
\[
64^{L+1}s\equiv d-64p\pmod{81}.
\]
This is the stated formula.  Multiplication by the unit $64^{-(L+1)}$ preserves the distinctness of the three values of $d$.
\end{proof}

\section{The exact two-carry recursion}

Fix $u\in\E^k$ and $v\in\E^L$.  Put
\[
r=R_k(u),
\qquad
q=Q_L(v),
\qquad
A=64^k,
\qquad
M=81^L.
\]
Let $c\in[0,M)$ and $h\in\Z$ be defined by
\[
r+Ac=q+Mh.
\]
Thus $c=c_{k,L}(u,v)$, while $h$ is the complementary quotient of the same Chinese-remainder lift.

For extensions $ua$ and $dv$, write
\[
r'=R_{k+3}(ua),
\qquad
q'=Q_{L+1}(dv),
\]
and define $c'\in[0,81M)$ and $h'\in\Z$ by
\[
r'+ABc'=q'+81Mh'.
\]

\begin{theorem}[Exact block carry recursion]
\label{thm:block-recursion}
Let
\[
t=t_k(u,a),
\qquad
s=s_L(v,d).
\]
There is a unique integer $z$ such that
\[
Bc'=c-t+Mz
\]
and
\[
81h'=h-s+Az.
\]
Conversely, these two equations are equivalent to the equality of the old and new Chinese-remainder lifts.
\end{theorem}

\begin{proof}
Lemma \ref{lem:lifts} gives
\[
r'=r+At,
\qquad
q'=q+Ms.
\]
Using
\[
r+Ac=q+Mh
\]
and
\[
r'+ABc'=q'+81Mh',
\]
we obtain
\begin{align*}
A(Bc'-c+t)
&=M(81h'+s-h).
\end{align*}
Since
\[
\gcd(A,M)=1,
\]
there is a unique integer $z$ satisfying
\[
Bc'-c+t=Mz
\]
and
\[
81h'+s-h=Az.
\]
Rearranging gives the displayed recursion.  Reversing these algebraic steps proves the converse.
\end{proof}

\begin{corollary}[Pinned one-block synchronization]
\label{cor:pinned-sync}
Assume $L\ge3$.  Then the extension pair $(ua,dv)$ has new carry
\[
c'=0
\]
if and only if
\[
c=t_k(u,a)
\]
and
\[
h\equiv s_L(v,d)\pmod{81}.
\]
For a fixed old pair $(u,v)$, at most one choice of
\[
(a,d)\in\E^3\times\E
\]
produces a pinned extension.
\end{corollary}

\begin{proof}
Here
\[
M=81^L>64^3=B.
\]
If $c'=0$, the first equation of Theorem \ref{thm:block-recursion} gives
\[
c-t=-Mz.
\]
Because
\[
-(B-1)\le c-t\le M-1
\]
and $M>B$, the only multiple of $M$ in this interval is $0$.  Hence
\[
c=t,
\qquad
z=0.
\]
The second equation then says
\[
h-s=81h',
\]
which is equivalent to the stated congruence.  The converse follows by setting
\[
z=0,
\qquad
h'=(h-s)/81.
\]
Distinctness of the local lift digits in Lemma \ref{lem:lifts} proves uniqueness of $a$ and $d$.
\end{proof}

Let
\[
\Dset_G=\{0,19,62\}
\]
be the domain of $G$ modulo $81$.

\begin{theorem}[Dynamical form of critical renormalization]
\label{thm:dynamical-renorm}
Assume $L\ge3$.  Pinned pairs at parameters $(k+3,L+1)$ are in bijection with old pairs $(u,v)\in\E^k\times\E^L$ satisfying all three conditions
\[
0\le c_{k,L}(u,v)<B,
\]
\[
Q_k(u)+81^k c_{k,L}(u,v)\bmod B\in\Rset_3,
\]
and
\[
R_L(v)+64^L h_{k,L}(u,v)\bmod81\in\Dset_G,
\]
where $h_{k,L}(u,v)$ is defined by
\[
R_k(u)+64^k c_{k,L}(u,v)
=Q_L(v)+81^L h_{k,L}(u,v).
\]
The corresponding three-symbol block and one-symbol block are unique.
\end{theorem}

\begin{proof}
By Corollary \ref{cor:pinned-sync}, a pinned extension requires
\[
c=t_k(u,a).
\]
Since $0\le t_k(u,a)<B$, this implies $c<B$.  The definition of $t_k(u,a)$ shows that such an $a$ exists precisely when
\[
Q_k(u)+81^kc\bmod B\in\Rset_3.
\]
Injectivity of the realization map makes $a$ unique.

Similarly, the congruence
\[
h\equiv s_L(v,d)\pmod{81}
\]
and the formula in Lemma \ref{lem:lifts} are equivalent to
\[
64R_L(v)+64^{L+1}h\equiv d\pmod{81}.
\]
Put
\[
y=R_L(v)+64^Lh.
\]
The last congruence says that $64y-d$ is divisible by $81$.  There is a symbol $d\in\E$ with this property precisely when
\[
y\bmod81\in\{0,19,62\}=\Dset_G.
\]
The residue determines $d$ uniquely.  Corollary \ref{cor:pinned-sync} now gives both directions of the bijection.
\end{proof}

For $0\le c<B$, $\xi\in\Z/B\Z$, and $\eta\in\Z/81\Z$, define the decorated low-carry profile
\[
P_{k,L}(c,\xi,\eta)
\]
to be the number of pairs $(u,v)\in\E^k\times\E^L$ for which
\[
c_{k,L}(u,v)=c,
\]
\[
Q_k(u)+81^kc\equiv\xi\pmod B,
\]
and
\[
R_L(v)+64^Lh_{k,L}(u,v)\equiv\eta\pmod{81}.
\]

\begin{corollary}[Exact decorated-profile formula]
\label{cor:profile}
For $L\ge3$,
\[
C_{k+3,L+1}(0)
=
\sum_{c=0}^{B-1}
\sum_{\xi\in\Rset_3}
\sum_{\eta\in\Dset_G}
P_{k,L}(c,\xi,\eta).
\]
Consequently,
\[
C_{k+3,L+1}(0)
\le
\sum_{c=0}^{B-1}C_{k,L}(c),
\]
and
\[
N(k+3,L+1)
\le
\sum_{c=0}^{B-1}C_{k,L}(c)-1.
\]
\end{corollary}

\begin{proof}
The equality is Theorem \ref{thm:dynamical-renorm} with the accepted old pairs grouped by their three displayed residues.  Dropping the two local acceptance restrictions gives the first inequality.  The all-zero pair occurs once in the low-carry sum and once in the new pinned fibre.  Using
\[
C_{k+3,L+1}(0)=N(k+3,L+1)+1
\]
gives the second inequality.
\end{proof}

\begin{proposition}[The local masks have exact sizes]
\label{prop:mask-sizes}
For a fixed right word $u\in\E^k$, exactly $27$ values
\[
c\in[0,B)
\]
satisfy
\[
Q_k(u)+81^kc\bmod B\in\Rset_3.
\]
For a fixed left word $v\in\E^L$, exactly three residue classes
\[
h\pmod{81}
\]
satisfy
\[
R_L(v)+64^Lh\bmod81\in\Dset_G.
\]
The first local density is therefore
\[
\frac{27}{64^3},
\]
and the second is
\[
\frac{3}{81}=\frac1{27}.
\]
\end{proposition}

\begin{proof}
Multiplication by $81^k$ is a permutation of $\Z/B\Z$.  Hence
\[
c\longmapsto Q_k(u)+81^kc\pmod B
\]
is a bijection.  The set $\Rset_3$ has $3^3=27$ elements, giving the first count.

Multiplication by $64^L$ is a permutation of $\Z/81\Z$.  Hence
\[
h\longmapsto R_L(v)+64^Lh\pmod{81}
\]
is a bijection.  The set $\Dset_G$ has three elements, giving the second count.
\end{proof}

\begin{remark}
The two exact mask sizes are marginal statements.  They do not specify how the actual carry coordinates $(c,h)$ are distributed inside the product of the two masks, so they do not by themselves bound the joint sum in Corollary \ref{cor:profile}.
\end{remark}

\begin{proposition}[Periodicity of the local acceptance rule]
\label{prop:periods}
As a function of the decorated residue state, the right acceptance rule depends on $k$ only modulo $2^{14}$, and the left acceptance rule depends on $L$ only modulo $9$.
\end{proposition}

\begin{proof}
We first prove
\[
81^{2^{14}}\equiv1\pmod{2^{18}}.
\]
For $j\ge0$, put
\[
x_j=81^{2^j}.
\]
Since
\[
\nu_2(x_0-1)=\nu_2(80)=4
\]
and $x_j\equiv1\pmod8$, one has
\[
\nu_2(x_j+1)=1.
\]
Therefore
\[
\nu_2(x_{j+1}-1)
=
\nu_2(x_j-1)+\nu_2(x_j+1)
=
\nu_2(x_j-1)+1.
\]
Induction gives
\[
\nu_2\bigl(81^{2^j}-1\bigr)=j+4.
\]
At $j=14$ the valuation is $18$, so multiplication by $81^k$ modulo
\[
B=2^{18}
\]
is periodic in $k$ with period $2^{14}$.

The successive powers of $64$ modulo $81$ are
\[
64,46,28,10,73,55,37,19,1.
\]
Thus
\[
64^9\equiv1\pmod{81}.
\]
The two acceptance conditions in Theorem \ref{thm:dynamical-renorm} use only $81^k$ modulo $B$ and $64^L$ modulo $81$, proving the assertion.
\end{proof}

\section{Multi-block synchronization before the first wrap}

Let $k,L,J\ge1$.  Begin with words
\[
u_0\in\E^k,
\qquad
v_0\in\E^L.
\]
For $1\le i\le J$, choose
\[
a_i\in\E^3,
\qquad
d_i\in\E,
\]
and define
\[
u_i=u_{i-1}a_i,
\qquad
v_i=d_iv_{i-1}.
\]
Let $c_i,h_i$ be the carry coordinates of $(u_i,v_i)$ at parameters
\[
(k+3i,L+i).
\]
Let
\[
t_i=t_{k+3(i-1)}(u_{i-1},a_i),
\qquad
s_i=s_{L+i-1}(v_{i-1},d_i).
\]

\begin{theorem}[No-wrap synchronization over several critical blocks]
\label{thm:nowrap}
Assume that the top pair is pinned,
\[
c_J=0,
\]
and that
\[
B^J<81^L.
\]
Then every intermediate wrap integer in Theorem \ref{thm:block-recursion} is zero.  More precisely,
\[
c_{i-1}=t_i+Bc_i
\]
and
\[
h_{i-1}=s_i+81h_i
\]
for every $1\le i\le J$.  Consequently,
\[
c_0=t_1+Bt_2+\cdots+B^{J-1}t_J
\]
and
\[
h_0=s_1+81s_2+\cdots+81^{J-1}s_J+81^Jh_J.
\]
In particular, the $J$ forward lift digits are recovered uniquely from the base-$B$ expansion of $c_0$, and the $J$ backward lift digits are recovered uniquely from $h_0$ modulo $81^J$.
\end{theorem}

\begin{proof}
At the step from level $i-1$ to level $i$, Theorem \ref{thm:block-recursion} gives an integer $z_i$ such that
\[
Bc_i=c_{i-1}-t_i+81^{L+i-1}z_i.
\]
Equivalently,
\[
c_{i-1}=t_i+Bc_i-81^{L+i-1}z_i.
\]
We descend from $i=J$ to $i=1$.  We claim that after $J-i$ descents,
\[
0\le c_i\le B^{J-i}-1.
\]
At $i=J$ this is true because $c_J=0$.

Assume it holds at level $i$.  Then
\begin{align*}
0
&\le t_i+Bc_i\\
&\le (B-1)+B(B^{J-i}-1)\\
&=B^{J-i+1}-1\\
&\le B^J-1\\
&<81^L\\
&\le81^{L+i-1}.
\end{align*}
Both $c_{i-1}$ and the integer $t_i+Bc_i$ therefore lie in
\[
[0,81^{L+i-1}).
\]
Their difference is a multiple of $81^{L+i-1}$, so they are equal.  Thus
\[
z_i=0,
\qquad
c_{i-1}=t_i+Bc_i,
\]
and the claimed bound continues to level $i-1$.

The second equation of Theorem \ref{thm:block-recursion}, with $z_i=0$, gives
\[
h_{i-1}=s_i+81h_i.
\]
Iterating the two recurrences yields the displayed expansions.  Since
\[
0\le t_i<B
\]
and
\[
0\le s_i<81,
\]
ordinary uniqueness of positional expansions recovers the digits as stated.
\end{proof}

\begin{corollary}[A positive fraction of every long pinned trace is synchronized]
\label{cor:fraction}
Let a pinned pair have parameters $(K,L_{\rm top})$.  Let $J$ satisfy
\[
0\le J\le\min\left(\left\lfloor\frac K3\right\rfloor,L_{\rm top}\right)
\]
and
\[
(64^3)^J<81^{L_{\rm top}-J}.
\]
Then it may be stripped through $J$ critical blocks without any carry wrap.  Consequently, for every real number
\[
0<\theta<\frac{\log81}{\log(81\cdot64^3)},
\]
and all sufficiently large $L_{\rm top}$, the first
\[
J=\lfloor\theta L_{\rm top}\rfloor
\]
stripped blocks synchronize as in Theorem \ref{thm:nowrap}.  The limiting fraction is
\[
\frac{\log81}{\log(81\cdot64^3)}=0.260472\ldots.
\]
\end{corollary}

\begin{proof}
After stripping $J$ blocks, the remaining backward length is
\[
L=L_{\rm top}-J.
\]
The first assertion is Theorem \ref{thm:nowrap} with this value of $L$.

The inequality
\[
(64^3)^J<81^{L_{\rm top}-J}
\]
is equivalent to
\[
J\log(64^3)<(L_{\rm top}-J)\log81,
\]
or
\[
\frac{J}{L_{\rm top}}
<
\frac{\log81}{\log81+\log(64^3)}.
\]
If $\theta$ is strictly below the right side, then
\[
J=\lfloor\theta L_{\rm top}\rfloor
\]
satisfies the inequality for every sufficiently large $L_{\rm top}$.  The decimal is the numerical value of the displayed exact constant.
\end{proof}

\section{An explicit off-zero source for a pinned extension}

The next calculation shows that the old zero carry is not an inductively closed state.

\begin{proposition}[A pinned extension arising from carry $257792$]
\label{prop:offzero}
Let
\[
\begin{aligned}
u={}&(1,0,1,0,1,0,1,1,0,0,-1,0),\\
v={}&(-1,0,-1,1),\\
a={}&(0,-1,-1),\\
d={}&1.
\end{aligned}
\]
For the old pair $(u,v)$ at parameters $(12,4)$, the carry coordinates are
\[
c=257792
\]
and
\[
h=28280728077103985211.
\]
This carry is nonzero and satisfies
\[
c<64^3.
\]
The pair passes both local tests in Theorem \ref{thm:dynamical-renorm}; its unique extension $(ua,dv)$ is pinned at parameters $(15,5)$.  Its canonical start is
\[
Y=1217392611211961739383252815.
\]
\end{proposition}

\begin{proof}
The word numerators and canonical data are
\[
N_{12}(u)=2318912255527500481137,
\]
\[
R_{12}(u)=4310860030160454133583,
\]
\[
Q_{12}(u)=72815604730280023521600,
\]
\[
N_4(v)=-601073,
\]
\[
R_4(v)=6707377,
\qquad
Q_4(v)=17209684,
\]
and
\[
N_3(a)=-9280,
\qquad
R_3(a)=239680,
\qquad
Q_3(a)=485900.
\]
These values are certified by the exact identities
\[
64^{12}Q_{12}(u)-81^{12}R_{12}(u)=N_{12}(u),
\]
\[
64^4Q_4(v)-81^4R_4(v)=N_4(v),
\]
and
\[
64^3Q_3(a)-81^3R_3(a)=N_3(a),
\]
together with the defining ranges of the canonical representatives.

Put
\[
A=64^{12}=4722366482869645213696,
\]
\[
M=81^4=43046721,
\qquad
B=64^3=262144.
\]
Direct integer multiplication gives
\begin{align*}
R_{12}(u)+A\cdot257792
&=1217392611211961739383252815\\
&=Q_4(v)+M\cdot28280728077103985211.
\end{align*}
Thus the old carry coordinates are the stated $c$ and $h$.  Also
\[
0<257792<262144=B.
\]

For the right local test,
\[
Q_{12}(u)+81^{12}c
\equiv239680
\equiv R_3(a)
\pmod B.
\]
Hence
\[
t_{12}(u,a)=c.
\]
For the left local test,
\[
R_4(v)+64^4h\equiv19\pmod{81}.
\]
Since $19$ is the residue corresponding to the symbol $d=1$, the backward test is satisfied.  More explicitly,
\[
h\equiv9\pmod{81},
\]
so Lemma \ref{lem:lifts} gives
\[
s_4(v,1)=9
\]
and
\[
Q_5(dv)=Q_4(v)+9M=404630173.
\]
Theorem \ref{thm:dynamical-renorm} now gives the pinned extension.  Its right canonical start is
\[
R_{15}(ua)=R_{12}(u)+Ac=Y.
\]
Finally,
\[
Y-Q_5(dv)
=81^5\cdot349144791075357842,
\]
so the pinning congruence modulo $81^5$ is also visible directly.
\end{proof}

\begin{corollary}[The old zero fibre is not a closed induction state]
\label{cor:notclosed}
The exact recursion for $C_{k+3,L+1}(0)$ cannot be obtained by retaining only old pairs of carry zero.  In particular, the contribution in Proposition \ref{prop:offzero} would be lost by such a recursion.
\end{corollary}

\begin{proof}
The old pair in Proposition \ref{prop:offzero} has carry
\[
c=257792\ne0,
\]
but its extension has new carry zero.  Therefore off-zero old carry fibres can contribute to the new pinned fibre.
\end{proof}

\section{Assessment}

This round produces a genuine one-block synchronization theorem.  The pinned fibre at $(k+3,L+1)$ is not an opaque intersection: for $L\ge3$ it is exactly the accepted part of the fixed old carry window
\[
0\le c<64^3,
\]
with one explicit three-step forward mask and one explicit one-step backward mask.  Every old pair has at most one pinned extension.  Corollary \ref{cor:profile} is therefore an exact critical renormalization, not a heuristic transfer.

The multi-block theorem goes further.  The local injections multiply without loss until the first modular wrap.  For any pinned trace, a positive fraction approaching
\[
0.260472\ldots
\]
of its critical blocks is synchronized by ordinary positional expansions.  This is the strongest unconditional synchronization statement obtained so far in this branch.

The result does not close the zero-trace conjecture.  Proposition \ref{prop:offzero} shows why the scalar target
\[
C_{k,L}(0)=1
\]
is not inductively closed: a new pinned solution may come from a nonzero old carry near the top of the fixed $64^3$ window.  After enough blocks are stripped, the inequality in Theorem \ref{thm:nowrap} fails, a nonzero wrap integer $z$ appears, and the base-$64^3$/base-$81$ digit synchronization can change discontinuously.  The first wrap, rather than the all-carry average or global carry injectivity, is now the exact live obstruction.

No nontrivial Collatz cycle and no divergent Collatz trajectory is excluded by this round.  Even a complete proof of pinned carry repulsion would still settle only the recurrent $81/64$ subsystem; the global synchronization step connecting every hypothetical least Collatz counterexample to that subsystem remains open.

\section{Next-round directive}

Continue on branch \texttt{P3-ZT-PINNED-CARRY}.  Use Theorem \ref{thm:block-recursion} as the basic state evolution and treat the first nonzero wrap integer $z$ as the new target.  Do not launch the monolithic $(k,L)=(28,10)$ search yet: at the present cost it returns one bit after roughly $3^{22}$ streamed heads, while the computations below directly test the new proof mechanism.

First finish and record the count-all runs for
\[
(23,8),(24,7),(24,8),(25,8),(25,9),
\]
including $N(k,L)$, $C_{k,L}(0)$, all recovered witnesses when the count is at most $1000$, split $h$, source and output SHA-256 digests, Python version, peak resident memory, and elapsed wall-clock time.  Preserve the already completed counts for $(20,7)$ through $(23,7)$ in the same format.

Implement \texttt{compute/zt\_carry\_small.py} exactly as directed in Round 2.  Verify the full carry spectra for
\[
(4,2),(5,2),(6,2),(7,3),(8,3),(9,3),(10,3),
\]
the support size $5^L$, the signed multiplicity formula, the energy $19^L$, and the second-moment identity.  This is a cheap independent audit of every load-bearing identity from Round 2 and should be completed before deeper computation is booked.

Then implement \texttt{compute/zt\_strip\_profile.py}.  For each pinned solution at
\[
(15,5),(20,7),(21,7),(22,7),(23,7),(24,7),(25,8),
\]
recover the complete right word and left word, not only the midpoint integer.  Strip successive $3$-symbol/$1$-symbol blocks.  At every level record
\[
(k_i,L_i,c_i,h_i,t_i,s_i,z_i),
\]
and verify the two equations
\[
64^3c_i=c_{i-1}-t_i+81^{L_{i-1}}z_i,
\]
\[
81h_i=h_{i-1}-s_i+64^{k_{i-1}}z_i.
\]
For each solution record the first index at which $z_i\ne0$, the value and sign of that first wrap, the residues of the right endpoint modulo $64^3$, and the backward ancestor modulo $81$.  Verify Theorem \ref{thm:nowrap} on every solution by confirming $z_i=0$ throughout the range
\[
(64^3)^J<81^{L_{\rm top}-J}.
\]

Next build a count-only version of the decorated profile in Corollary \ref{cor:profile}.  For each predecessor cell
\[
(12,4),(17,6),(18,6),(19,6),(20,6),(21,6),(22,7),
\]
count old pairs with $0\le c<64^3$ in the four classes determined by whether the right three-step test and the left one-step test pass.  Record
\[
T_{00},T_{01},T_{10},T_{11},
\]
where
\[
T_{11}=C_{k+3,L+1}(0).
\]
Also record the right and left marginals and the exact covariance
\[
T_{11}(T_{00}+T_{01}+T_{10}+T_{11})
-(T_{10}+T_{11})(T_{01}+T_{11}).
\]
Validate this profile by complete enumeration for every feasible pair with $k\le10$ and $L\le3$ before using meet-in-the-middle output.

The theoretical task next round is to classify the first wrap.  Starting from Theorem \ref{thm:block-recursion}, prove the strongest true restriction on $z$ at the first index where it is nonzero.  The preferred target is a theorem that, in the critical region $k\le3L-2$, the first wrap cannot simultaneously satisfy the right set $\Rset_3$ and the left domain $\{0,19,62\}$.  If that statement is false, return the smallest exact counterexample with its complete word data and replace it by the strongest residue restriction surviving all computed wrap signatures.

Edit \texttt{COLLATZ\_MAP.md} to record as proved in Round 3: the exact two-carry recursion, pinned one-block synchronization, the decorated-profile formula, the low-carry inequality, the periodic local masks, and the multi-block no-wrap theorem.  Record the explicit carry-$257792$ extension as the reason the old zero fibre is not an inductively closed state.  Rename the immediate obstruction from \texttt{pinned carry repulsion} to \texttt{first-wrap exclusion}, while retaining pinned carry repulsion as the branch's final finite target.

\end{document}
